\begin{definition}[The translated swept-string patch]%
    \label{def:peps_torus_swept_string_patch}
    \lean{TNLean.PEPS.torusSweptStringPatch,
        TNLean.PEPS.torusSweptStringVertices,
        TNLean.PEPS.torusSweptStringBottomVertices}
    \leanok
    Let the torus periods satisfy $w\geq5$ and $H\geq4$.
    For $v=(x,y)$, with all coordinates read periodically, define
    \begin{align}
        P_v&=\{x,x+1,x+2,x+3\}\times\{y,y+1,y+2\},\\
        S_1&=\{x+1,x+2\}\times\{y+1\},\\
        B&=\{x+1,x+2\}\times\{y\},& S_2&=S_1\cup B.
    \end{align}
    These finite sets specify the two successive sweeps in the virtual
    string-deformation figure of \cite[Section~6.6]{Schuch2010PEPS},
    local source lines 2361--2386.
    % Source label: eq:anyons:fluxon-braiding-lazy.
\end{definition}

\begin{definition}[Literal initial and swept string assignments]%
    \label{def:peps_torus_swept_string_assignment}
    \lean{TNLean.PEPS.torusSweptStringInitialRight,
        TNLean.PEPS.torusSweptStringInitialUp,
        TNLean.PEPS.torusSweptStringInitialOperators,
        TNLean.PEPS.torusSweptStringOperators}
    \leanok
    \uses{def:peps_torus_swept_string_patch,def:peps_regular_swept_patch}
    For any group $G$ and $g,h\in G$, prescribe the initial rightward
    and upward transports by
    \begin{align}
        D_{\rightarrow}(p)&=
        \begin{cases}
            h^{-1},&p_x=x+1, p_y\in\{y-1,y,y+1\},\\
            1,&\text{otherwise},
        \end{cases}\\
        D_{\uparrow}(p)&=
        \begin{cases}
            g,&p_x\in\{x,x+1,x+2,x+3\}, p_y=y+1,\\
            1,&\text{otherwise}.
        \end{cases}
        \label{eq:peps_swept_string_initial}
    \end{align}
    Assign each native right or up edge the prescribed transport if
    its forward direction agrees with the ordered orientation, and its
    inverse otherwise. Denote this ordered assignment by $u^{(0)}$,
    and put $u^{(j)}=(u^{(0)})^{S_j,g}$ for $j=1,2$.
    The second string continues below the displayed patch to retain
    its distant partner. A displayed rightward or downward matrix
    labelled $s$ has transport $s^{-1}$ in the corresponding direction.
    This is a finite realization of the virtual figure in
    \cite[Section~6.6]{Schuch2010PEPS}, lines 2361--2386.
\end{definition}

\begin{theorem}[Native transport tables, including both seams]%
    \label{thm:peps_torus_swept_string_transport}
    \lean{TNLean.PEPS.torusSweptStringInitialOperators_right_transport,
        TNLean.PEPS.torusSweptStringInitialOperators_up_transport,
        TNLean.PEPS.torusSweptStringOperators_right_transport,
        TNLean.PEPS.torusSweptStringOperators_up_transport,
        TNLean.PEPS.torusSweptStringOperators_right_ordered,
        TNLean.PEPS.torusSweptStringOperators_up_ordered}
    \leanok
    \uses{def:peps_torus_swept_string_assignment,
        thm:peps_regular_swept_patch_transport}
    The actual initial rightward and upward transports are exactly
    (\ref{eq:peps_swept_string_initial}). Each swept directed transport
    is given by (\ref{eq:peps_swept_directed_cases}), with $R=S_j$ and
    the corresponding initial transport.
    Write $\bar z$ for the least nonnegative representative of a cyclic
    coordinate $z$. If $T$ is this swept rightward transport, its ordered
    edge coefficient is $T$ when $\overline{p_x}<\overline{p_x+1}$ and
    $T^{-1}$ otherwise. The analogous upward coefficient uses
    $\overline{p_y}<\overline{p_y+1}$.
    These formulas hold on every native edge, including crossing bonds
    of $P_v$ and edges on either torus seam.
\end{theorem}

\begin{proof}\leanok
    \uses{thm:peps_regular_swept_patch_transport}
    Every native edge is uniquely a right or up edge. Its ordered
    orientation agrees with the forward direction precisely when the
    stated comparison of representatives holds. The two inverse
    corrections cancel in its initial directed transport.
    Apply the four directed gauge cases, then invert again when
    recovering the ordered edge coefficient.
\end{proof}

\begin{theorem}[The second sweep acts only on the newly added vertices]%
    \label{thm:peps_torus_swept_string_successive}
    \lean{TNLean.PEPS.torusSweptStringOperators_second_eq_sweep}
    \leanok
    \uses{def:peps_torus_swept_string_assignment}
    The literal assignments satisfy $u^{(2)}=(u^{(1)})^{B,g}$.
    Thus no vertex in $S_1$ receives the gauge $g^{-1}$ twice.
\end{theorem}

\begin{proof}\leanok
    The sets $S_1$ and $B$ are disjoint and have union $S_2$.
    At each endpoint, composing the two gauges therefore gives
    $g^{-1}$ on the union and identity elsewhere. Substitution into
    the ordered endpoint formula proves the asserted equality.
\end{proof}

\begin{theorem}[Actual closed coefficient equality of the three presentations]%
    \label{thm:peps_torus_swept_string_state}
    \lean{TNLean.PEPS.stateCoeff_torusSweptStringOperators,
        TNLean.PEPS.stateCoeff_torusSweptStringOperators_first_eq_second}
    \leanok
    \uses{def:peps_torus_swept_string_assignment,
        thm:peps_regular_swept_patch_state}
    Suppose $G$ is finite and the incident site coefficients are
    invariant under simultaneous left translation. The actual closed
    coefficient vectors obey
    $\Psi_a(u^{(0)})=\Psi_a(u^{(1)})=\Psi_a(u^{(2)})$.
    The changed crossing bonds are included in each contraction;
    no equality of fixed boundary columns is asserted.
\end{theorem}

\begin{proof}\leanok
    \uses{thm:peps_regular_swept_patch_state}
    Apply the closed swept-patch identity separately to $S_1$ and $S_2$.
    Their equalities with the initial vector also equate the two swept vectors.
\end{proof}
